% Source: https://artofproblemsolving.com/community/c3h1356161_similar_triangles_in_parallelograms
% Stable identity: aops-topic-1356161-q01
\begin{asy}
usepackage("amsmath");
size(5cm);
defaultpen(fontsize(10pt));
pair A=(0,0),B=(6,0),C=(8,4),D=(2,4),Em=(3,0);
pair F=extension(A,C,Em,D);
fill(A--F--Em--cycle,paleblue); fill(C--F--D--cycle,palered);
draw(A--B--C--D--cycle,linewidth(0.9));
draw(A--C); draw(Em--D);
dot(Em);dot(F);
label("$A$",A,SW); label("$B$",B,SE); label("$C$",C,NE); label("$D$",D,NW); label("$E$",Em,S); label("$F$",F,SE);
label("\mbox{}",(min(currentpicture,user=true)+max(currentpicture,user=true))/2,black); // keeps the text after the figure black when LaTeX typesets the labels inline
\end{asy}

The angles $\angle AFE$ and $\angle CFD$ are vertical angles, so they are equal. Since $AB\parallel CD$ and $AC$ is a transversal, the alternate angles $\angle FAE=\angle CAB$ and $\angle FCD=\angle ACD$ are equal. By AA, $\triangle AFE\sim\triangle CFD$ (with $A\leftrightarrow C$, $F\leftrightarrow F$, $E\leftrightarrow D$). $\blacksquare$

The similarity ratio is $\dfrac{AE}{CD}=\dfrac{\dfrac{1}{2}AB}{AB}=\dfrac{1}{2}$, so $AF:FC=\boxed{1:2}$; that is, $F$ lies one third of the way from $A$ to $C$.
